\documentclass[11pt]{article}

\usepackage[T1]{fontenc}
\usepackage[utf8]{inputenc}
\usepackage{lmodern}
\usepackage[margin=1.04in]{geometry}
\usepackage{amsmath,amssymb,mathtools,amsthm}
\usepackage{aliascnt}
\usepackage{mathrsfs}
\usepackage{microtype}
\usepackage{enumitem}
\usepackage{xcolor}
\usepackage{url}
\usepackage{hyperref}
\usepackage[capitalize,noabbrev]{cleveref}

\hypersetup{
  colorlinks=true,
  linkcolor=blue!45!black,
  citecolor=green!35!black,
  urlcolor=blue!55!black,
  pdftitle={A Polynomial Kontsevich--Zagier Theorem for Compact Exponential Periods on the Affine Line},
  pdfauthor={Christopher D. Long; Antoine-Auguste Le Blanc},
  pdfsubject={Exponential periods, polynomial relations, E-functions, formal periods},
  pdfkeywords={exponential period, formal period, Kontsevich--Zagier conjecture, E-function, polynomial relations, affine line, differential Galois theory}
}

\numberwithin{equation}{section}
\setlength{\emergencystretch}{3em}

\newtheorem{theorem}{Theorem}[section]
\newaliascnt{proposition}{theorem}
\newtheorem{proposition}[proposition]{Proposition}
\aliascntresetthe{proposition}
\newaliascnt{lemma}{theorem}
\newtheorem{lemma}[lemma]{Lemma}
\aliascntresetthe{lemma}
\newaliascnt{corollary}{theorem}
\newtheorem{corollary}[corollary]{Corollary}
\aliascntresetthe{corollary}
\theoremstyle{definition}
\newaliascnt{definition}{theorem}
\newtheorem{definition}[definition]{Definition}
\aliascntresetthe{definition}
\theoremstyle{remark}
\newaliascnt{remark}{theorem}
\newtheorem{remark}[remark]{Remark}
\aliascntresetthe{remark}

\newcommand{\C}{\mathbb C}
\newcommand{\Qbar}{\overline{\mathbb Q}}
\newcommand{\Aone}{\mathbb A^1}
\newcommand{\Spec}{\operatorname{Spec}}
\newcommand{\Span}{\operatorname{span}}
\newcommand{\Frac}{\operatorname{Frac}}
\newcommand{\Sym}{\operatorname{Sym}}
\newcommand{\per}{\operatorname{per}}
\newcommand{\ord}{\operatorname{ord}}
\newcommand{\Pcal}{\mathcal P}
\newcommand{\Acal}{\mathcal A}
\newcommand{\Ocal}{\mathcal O}
\newcommand{\dd}{\,d}

\title{A Polynomial Kontsevich--Zagier Theorem for\\
Compact Exponential Periods on the Affine Line}
\author{Christopher D. Long \and Antoine-Auguste Le Blanc\thanks{Le Blanc cryptographic author identifier: \texttt{b6048125\allowbreak{}30b3535b\allowbreak{}23d6dbb0\allowbreak{}f5abe32d\allowbreak{}1ed1a9de\allowbreak{}35fb4f7d\allowbreak{}04d62cff\allowbreak{}08ea19c8}}}
\date{August 2026}

\begin{document}
\maketitle

\begin{center}
\small
\texttt{galizur@gmail.com}
\end{center}

\begin{abstract}
Let $q,r\in\Qbar[x]$, let $\Delta$ be a finite degree-zero zero-cycle on
$\Aone(\Qbar)$ with coefficients in $\Qbar$, and let $\Gamma$ be any finite
piecewise $C^1$ one-chain in $\C$ with $\partial\Gamma=\Delta$.  We consider
compact exponential periods
\[
  \int_\Gamma r(x)e^{q(x)}\,dx
\]
and the zero-dimensional exponential periods $e^\lambda$, $\lambda\in\Qbar$.
The companion linear theorem identifies every numerical linear relation among
these periods with a formal consequence of bilinearity, functoriality, and
twisted Stokes.

We prove the corresponding polynomial theorem for the commutative algebra
generated by the same periods.  Every polynomial numerical relation is
generated by the companion linear relations, multiplication, and the
identities
\[
  e^0=1,\qquad e^\lambda e^\mu=e^{\lambda+\mu},\qquad
  e^\lambda\int_\Gamma r e^q\,dx
  =\int_\Gamma r e^{q+\lambda}\,dx.
\]
Equivalently, the period map is injective on the tensor-generated formal
algebra of compact exponential periods on the affine line.

The new input is a uniform description of the functional polynomial-relation
ideal.  After adjoining the finitely many exponential functions coming from
endpoint values and constant shifts, the affine Picard--Vessiot group of the
inhomogeneous moment system has a vector translation kernel and reductive
homogeneous quotient.  Horizontal projectors are Laurent polynomials in the
adjoined exponentials.  A Laurent-polynomial gauge separates the coordinates
already in the base ring from a full translation block, and the latter
coordinates are algebraically independent.  Hence the complete functional
relation ideal is generated by affine-linear relations.  Beukers' lifting
theorem transfers this statement to values, and the companion linear
Kontsevich--Zagier theorem makes every specialized generator formal.
\end{abstract}

\medskip
\noindent\textit{2020 Mathematics Subject Classification.}
Primary 11J85; Secondary 11J81, 11J91, 12H05, 14F10, 34M35.

\noindent\textit{Keywords.}
Exponential periods, formal periods, Kontsevich--Zagier period conjecture,
$E$-functions, polynomial relations, affine line, differential Galois theory.

\section{Introduction and main theorem}\label{sec:introduction}

Kontsevich and Zagier formulated their period conjecture as the assertion that
all identities between integral presentations of ordinary periods follow from
additivity, algebraic change of variables, and Stokes' formula
\cite{KontsevichZagier}.  In cohomological language this is the injectivity of
a formal period map; see \cite[Chapters~12--13]{HuberMullerStach} and
\cite{CressonViuSos}.  Exponential periods replace an algebraic differential
form $\omega$ by $e^f\omega$, where $f$ is a regular function, and replace the
de Rham differential by
\[
  d_f=d+df\wedge.
\]
General definitions and comparison results for exponential periods are
developed in \cite{CommelinHabeggerHuber}.  We use \emph{exponential period
function} in the sense of Fres\'an and Jossen \cite{FresanJossen}.

This paper is the nonlinear continuation of two companion papers.  The first
proves algebraic-value rigidity for
\[
  I_{q,\Delta}(z)=\int_\Gamma e^{zq(x)}\,dx
\]
and supplies the homogeneous moment connection and its residue structure
\cite{LongEIT}.  The second treats arbitrary polynomial differentials,
different polynomials in the same relation, and zero-dimensional exponential
periods, and proves the linear Kontsevich--Zagier theorem for compact
exponential periods on the affine line \cite{LongLinearKZ}.  The present paper
adds multiplication and classifies all polynomial relations in the resulting
tensor-generated algebra.

Throughout, put $k=\Qbar$.  For $q,r\in k[x]$ and a finite degree-zero
zero-cycle
\[
  \Delta=\sum_{\alpha}d_\alpha[\alpha]
  \quad\text{on }\Aone(k),
\]
choose a finite piecewise $C^1$ one-chain $\Gamma$ in $\C$ with
$\partial\Gamma=\Delta$ and set
\begin{equation}\label{eq:I}
  I_{q,r,\Delta}(z)=\int_\Gamma r(x)e^{zq(x)}\,dx.
\end{equation}
The one-form in \eqref{eq:I} has an entire primitive in $x$, so the function
is independent of the chosen chain and is entire in $z$.  The compact
exponential period is $I_{q,r,\Delta}(1)$.  A common nonzero algebraic
evaluation parameter is already absorbed into the polynomial $q$.

Let $\Pcal_{\Aone}^{\mathrm{KZ,cpt}}$ be the formal linear period space from
\cite{LongLinearKZ}.  It is generated over $k$ by the symbols
\[
  \mathsf P(q,r,\Delta)
  \quad\text{and}\quad
  \mathsf E(\lambda),\qquad \lambda\in k,
\]
modulo bilinearity, functoriality through smooth affine curves, and twisted
Stokes.  The companion theorem states that its period map
\[
  \mathsf P(q,r,\Delta)\longmapsto I_{q,r,\Delta}(1),
  \qquad
  \mathsf E(\lambda)\longmapsto e^\lambda
\]
is injective.

\begin{definition}[Tensor-generated formal algebra]\label{def:formal-algebra}
Let $\Acal_{\Aone}^{\mathrm{KZ,cpt}}$ be the quotient of
$\Sym_k(\Pcal_{\Aone}^{\mathrm{KZ,cpt}})$ by the ideal generated by
\begin{align}
  \mathsf E(0)&=1,\label{eq:E-zero}\\
  \mathsf E(\lambda)\mathsf E(\mu)
    &=\mathsf E(\lambda+\mu),\label{eq:E-product}\\
  \mathsf E(\lambda)\mathsf P(q,r,\Delta)
    &=\mathsf P(q+\lambda,r,\Delta).\label{eq:constant-shift}
\end{align}
Multiplication is the formal tensor, or Fubini, product.  The numerical
assignments above define a homomorphism
\begin{equation}\label{eq:period-map}
  \per:\Acal_{\Aone}^{\mathrm{KZ,cpt}}\longrightarrow\C.
\end{equation}
\end{definition}

\begin{theorem}[Polynomial Kontsevich--Zagier theorem]\label{thm:main}
The period map \eqref{eq:period-map} is injective.  Equivalently, every
polynomial numerical relation among compact exponential periods on the affine
line and zero-dimensional exponential periods is generated by
\begin{enumerate}[label=\textup{(\roman*)},leftmargin=2.4em]
\item the linear formal relations of \cite{LongLinearKZ};
\item multiplication;
\item \eqref{eq:E-zero}--\eqref{eq:constant-shift}.
\end{enumerate}
\end{theorem}

The theorem is not an algorithmic relation theorem for one fixed tuple of
$E$-functions.  Fischler and Rivoal proved that, for any fixed finite tuple of
$E$-functions, the ideals of functional and value relations can be computed
algorithmically \cite{FischlerRivoalEffective}; the exceptional-value and
minimal-operator problems were developed by Adamczewski--Rivoal and
Bostan--Rivoal--Salvy \cite{AdamczewskiRivoal,BostanRivoalSalvy}.  The new
point here is uniform and structural: for the entire geometric family
\eqref{eq:I}, after adjoining the required exponential shifts, the functional
ideal is generated by its linear part, and the companion paper identifies the
specialized generators by formal period operations.

The proof has three steps.  First, the moment connections from
\cite{LongEIT} remain semisimple after adjoining finitely many exponential
functions, and every horizontal projector over that exponential field is a
Laurent polynomial in those functions.  Second, affine Picard--Vessiot theory
splits the inhomogeneous moment coordinates into an affine-linear constrained
part and a full translation part.  The complete functional ideal is therefore
linear.  Third, Beukers' theorem lifts a numerical relation of arbitrary
degree to the functional ideal; specialization of its linear generators is
handled by \cite{LongLinearKZ}.

\section{Imported arithmetic and moment-theoretic input}
\label{sec:input}

\subsection{\texorpdfstring{$E$}{E}-functions and Beukers' lifting theorem}

We use Beukers' normalization of Siegel's $E$-functions
\cite[Section~1]{Beukers}.

\begin{definition}[$E$-function]
An entire function
\[
  F(z)=\sum_{n\ge0}a_n\frac{z^n}{n!}
\]
is an $E$-function if $a_n\in k$, the joint logarithmic height
$h(a_0,\ldots,a_n)$ is $O(n)$, and $F$ satisfies a nonzero linear
differential equation over $k(z)$.
\end{definition}

\begin{theorem}[Beukers]\label{thm:Beukers}
Let $F_1,\ldots,F_N$ be $E$-functions forming a vector solution of
\[
  \mathbf F'=A(z)\mathbf F,
  \qquad A(z)\in M_N(k(z)),
\]
and let $T\in k[z]$ clear the denominators of $A$.  Let
$\xi\in k$ satisfy $\xi T(\xi)\ne0$.  If a homogeneous polynomial
$P\in k[X_1,\ldots,X_N]$ satisfies
\[
  P(F_1(\xi),\ldots,F_N(\xi))=0,
\]
then there is a polynomial
$Q\in k[z,X_1,\ldots,X_N]$, homogeneous in the $X$-variables and of the
same degree as $P$, such that
\[
  Q(z,F_1(z),\ldots,F_N(z))\equiv0,
  \qquad
  Q(\xi,X_1,\ldots,X_N)=P(X_1,\ldots,X_N).
\]
\end{theorem}

This is the combination of \cite[Theorems~1.3 and~3.2]{Beukers}.  Adjoining
the constant $E$-function $1$ homogenizes an affine relation.

\begin{lemma}[Distinct exponential functions]\label{lem:distinct-exponentials}
If $\lambda_1,\ldots,\lambda_s\in\C$ are distinct, then
$e^{\lambda_1z},\ldots,e^{\lambda_sz}$ are linearly independent over
$\C(z)$.
\end{lemma}

\begin{proof}
Assume a nontrivial relation with the fewest nonzero terms.  Differentiate
and eliminate one term.  Minimality forces the logarithmic derivative of a
nonzero rational function to be a nonzero constant, which is impossible at
infinity.
\end{proof}

\subsection{The polynomial moment package}

Let $q\in k[x]$ have degree $n\ge2$, let $\Delta$ be a finite degree-zero
zero-cycle, and put
\[
  U_j(z)=\int_\Gamma x^j e^{zq(x)}\,dx,
  \qquad 0\le j\le n-2,
  \qquad
  \mathbf U=(U_0,\ldots,U_{n-2})^{\mathsf T}.
\]
The following is the exact form of the input proved in \cite{LongEIT} and
used in \cite{LongLinearKZ}.  No affine normalization of $q$ is required.

\begin{theorem}[Polynomial moment package]\label{thm:moment-package}
There are matrices $C_q,D_q\in M_{n-1}(k)$ and vectors
$\mathbf b_\lambda\in k^{n-1}$, indexed by
\[
  \Lambda=\{0\}\cup q(\operatorname{supp}\Delta),
\]
such that
\begin{equation}\label{eq:moment-system}
  \mathbf U'
  =\left(C_q+\frac{D_q}{z}\right)\mathbf U
   +\frac1z\sum_{\lambda\in\Lambda}
      \mathbf b_\lambda e^{\lambda z}.
\end{equation}
Every coordinate of $\mathbf U$ is an $E$-function.  The homogeneous module
\begin{equation}\label{eq:homogeneous}
  \mathbf Y'=\left(C_q+\frac{D_q}{z}\right)\mathbf Y
\end{equation}
is semisimple over $\C(z)$, and $C_q$ is semisimple.  Moreover,
\begin{equation}\label{eq:global-D-spectrum}
  \Spec(D_q)=
  \left\{-\frac1n,-\frac2n,\ldots,-\frac{n-1}{n}\right\}.
\end{equation}
If $\lambda$ is an eigenvalue of $C_q$, if $\pi_\lambda$ is the corresponding
spectral projector, and if
\[
  \overline D_{q,\lambda}
  =\pi_\lambda D_q|_{\ker(C_q-\lambda I)},
\]
then
\begin{equation}\label{eq:compressed-spectrum}
  \Spec(\overline D_{q,\lambda})
  =\bigcup_{\substack{\tau:q'(\tau)=0\\q(\tau)=\lambda}}
    \left\{-\frac{m_\tau}{m_\tau+1},
    -\frac{m_\tau-1}{m_\tau+1},\ldots,
    -\frac1{m_\tau+1}\right\},
\end{equation}
where $m_\tau=\ord_\tau q'$.
\end{theorem}

For a finite direct sum of such systems, we write
\begin{equation}\label{eq:direct-system}
  \mathbf U'=A\mathbf U+\mathbf f,
  \qquad
  A=C+\frac Dz,
  \qquad
  \mathbf f=\frac1z\sum_{\lambda\in\Lambda}
      \mathbf b_\lambda e^{\lambda z}.
\end{equation}
The matrices $C,D$ are block diagonal.  The spectra in
\eqref{eq:global-D-spectrum} and \eqref{eq:compressed-spectrum} therefore
remain contained in the open interval $(-1,0)$.

\section{Horizontal projectors over an exponential field}
\label{sec:projectors}

Let $\Gamma\subset k$ be a finitely generated additive subgroup containing
$\Lambda$.  Since $\Gamma$ is torsion-free, it is free abelian.  Put
\begin{equation}\label{eq:exponential-rings}
  A_\Gamma=\C[\Gamma],
  \qquad
  R_\Gamma=\C[z]\otimes_\C A_\Gamma,
  \qquad
  K_\Gamma=\Frac(R_\Gamma).
\end{equation}
The basis element $\varepsilon_\mu\in A_\Gamma$ is realized as
$e^{\mu z}$.  If $\omega_1,\ldots,\omega_r$ is a $\mathbb Z$-basis of
$\Gamma$, then
\[
  A_\Gamma=\C[e^{\pm\omega_1z},\ldots,e^{\pm\omega_rz}]
\]
is a Laurent polynomial ring.  By \cref{lem:distinct-exponentials}, this
realization is injective.  The differential constants of $K_\Gamma$ are
$\C$.  The torus $\mathbf T=(\C^\times)^r$ acts differentially by rescaling
the basis exponentials and fixes $\C(z)$.

\begin{lemma}[Finite torus-orbit decomposition]
\label{lem:finite-orbit}
Let $v\in K_\Gamma^m$.  If the complex span of its $\mathbf T$-orbit is
finite-dimensional, then
\[
  v=\sum_{\mu\in S}e^{\mu z}r_\mu(z)
\]
for a finite set $S\subset\Gamma$ and vectors
$r_\mu(z)\in\C(z)^m$.
\end{lemma}

\begin{proof}
The finite-dimensional orbit span is a rational representation of the torus
and hence is the direct sum of its character spaces.  A vector of weight
$\mu\in\Gamma$ becomes torus invariant after multiplication by
$e^{-\mu z}$, and therefore belongs to $\C(z)^m$.
\end{proof}

\begin{theorem}[Horizontal endomorphisms after exponential base change]
\label{thm:endomorphism-normal-form}
If $P\in M_N(K_\Gamma)$ satisfies
\begin{equation}\label{eq:horizontal-endomorphism}
  P'=[A,P],
  \qquad A=C+\frac Dz,
\end{equation}
then
\begin{equation}\label{eq:endomorphism-expansion}
  P(z)=\sum_{\mu\in S}e^{\mu z}P_\mu
\end{equation}
for a finite set $S\subset\Gamma$ and constant matrices
$P_\mu\in M_N(\C)$ satisfying
\begin{equation}\label{eq:intertwining}
  [C,P_\mu]=\mu P_\mu,
  \qquad
  [D,P_\mu]=0.
\end{equation}
In particular, every horizontal idempotent belongs to $M_N(A_\Gamma)$.
\end{theorem}

\begin{proof}
The solution space in $M_N(K_\Gamma)$ of
\eqref{eq:horizontal-endomorphism} has dimension at most $N^2$ over the
constant field and is stable under $\mathbf T$.  Thus the torus orbit of $P$
has finite-dimensional span.  By \cref{lem:finite-orbit},
\[
  P=\sum_{\mu\in S}e^{\mu z}R_\mu(z),
  \qquad R_\mu\in M_N(\C(z)).
\]
For each weight,
\begin{equation}\label{eq:Rmu-equation}
  R_\mu'+\mu R_\mu
  =[C,R_\mu]+\frac1z[D,R_\mu].
\end{equation}
A pole of $R_\mu$ at a nonzero finite point is impossible because
differentiation raises its order by one.  At zero, a leading term
$z^{-m}R_{-m}$ with $m\ge1$ would satisfy
\[
  -mR_{-m}=[D,R_{-m}].
\]
Every eigenvalue of $\operatorname{ad}D$ is a difference of two numbers in
$(-1,0)$, hence lies in $(-1,1)$; it cannot equal the negative integer
$-m$.  Thus $R_\mu$ is polynomial.

Suppose $R_\mu=R_dz^d+\cdots$ with $d\ge1$.  The leading coefficient gives
\[
  [C,R_d]=\mu R_d.
\]
Because $C$ is semisimple, a nonzero block of $R_d$ maps a
$C$-eigenspace $V_\alpha$ to $V_\beta$ with $\beta-\alpha=\mu$.
Projecting the next coefficient equation to this block gives
\[
  dR_{d,\beta\alpha}
  =\overline D_\beta R_{d,\beta\alpha}
   -R_{d,\beta\alpha}\overline D_\alpha.
\]
By \eqref{eq:compressed-spectrum}, the spectrum of the operator on the
right lies in $(-1,1)$, so it cannot contain the positive integer $d$.
Hence $d=0$, and \eqref{eq:Rmu-equation} yields
\eqref{eq:intertwining}.
\end{proof}

The appearance of exponential coefficients is necessary.  If two data differ
by a constant, then
\[
  \mathbf U_{q+c}(z)=e^{cz}\mathbf U_q(z),
\]
and the projector onto their graph has entries involving $e^{\pm cz}$.
Thus the correct rigidity statement is Laurent-polynomial, not constant,
after exponential base change.

\begin{lemma}[Semisimplicity after exponential base change]
\label{lem:base-change-semisimple}
The homogeneous module in \eqref{eq:direct-system} remains semisimple after
extension of scalars from $\C(z)$ to $K_\Gamma$.
\end{lemma}

\begin{proof}
Let $L_h/\C(z)$ be a Picard--Vessiot field for the homogeneous system and
let $G$ be its Galois group.  Semisimplicity of the faithful solution
representation implies that $G$ is reductive.  The extension
$K_\Gamma/\C(z)$ is Picard--Vessiot with torus Galois group.  In a common
differential overfield, put
\[
  J=L_h\cap K_\Gamma.
\]
The subgroup of the torus fixing $J$ is normal, so $J/\C(z)$ is itself a
Picard--Vessiot extension.  Therefore
$\operatorname{Gal}(L_h/J)$ is a closed normal subgroup of the reductive
group $G$, hence is reductive in characteristic zero.  The standard
base-change identification
\[
  \operatorname{Gal}(L_hK_\Gamma/K_\Gamma)
  \simeq\operatorname{Gal}(L_h/J)
\]
proves the assertion.  We use the Picard--Vessiot correspondence and
base-change results in \cite[Chapter~1]{vanDerPutSinger}.
\end{proof}

\section{Affine Picard--Vessiot linearity}
\label{sec:affine}

\begin{lemma}[Meromorphic homogeneous-solution exclusion]
\label{lem:meromorphic-exclusion}
Let
\[
  y'=\left(B(z)+\frac Dz\right)y,
\]
where $B$ is holomorphic at zero and
$\Spec(D)\cap\mathbb Z=\varnothing$.  A single-valued meromorphic solution
at zero is identically zero.  In particular, the operator
\[
  \mathscr L(v)=v'-Av
\]
is injective on $K_\Gamma^N$ for the direct-sum moment connection.
\end{lemma}

\begin{proof}
If a nonzero meromorphic solution has leading term
$z^m(v+O(z))$, where $m\in\mathbb Z$ and $v\ne0$, comparison of the
coefficient of $z^{m-1}$ gives $mv=Dv$.  This contradicts the hypothesis.
Every element of $K_\Gamma$ is single-valued and meromorphic at zero, and
\eqref{eq:global-D-spectrum} excludes integral residue eigenvalues.
\end{proof}

\begin{lemma}[Affine Picard--Vessiot exact sequence]
\label{lem:affine-exact-sequence}
Let $K$ be a differential field with algebraically closed constant field
$\C$.  Let $Y'=AY$ have reductive Picard--Vessiot group $H$ and constant
solution space $V$, and let $u'=Au+f$.  The Picard--Vessiot group
$\widetilde H$ of the augmented system fits into an exact sequence
\begin{equation}\label{eq:affine-exact-sequence}
  1\longrightarrow W\longrightarrow\widetilde H
  \longrightarrow H\longrightarrow1,
\end{equation}
where $W$ is an $H$-stable vector subspace of $V$.  If
$V=W\oplus V_0$ is $H$-stable, the translation term projected to $V_0$ is
a regular algebraic $1$-cocycle $H\to V_0$, and this cocycle is a
coboundary.
\end{lemma}

\begin{proof}
Choose a fundamental matrix
\[
  \widetilde Y=
  \begin{pmatrix}Y&u\\0&1\end{pmatrix}.
\]
Every $\sigma\in\widetilde H$ has a unique affine form
\[
  \sigma(Y)=Yh_\sigma,
  \qquad
  \sigma(u)=u+Yc_\sigma.
\]
The restriction map onto $H$ is surjective: an element of $K(Y)$ fixed by
its image is fixed by all of $\widetilde H$ and hence lies in $K$.  The
kernel is a closed algebraic subgroup of a vector group, therefore a vector
subspace $W$; conjugation makes it an $H$-submodule.

After choosing an $H$-stable complement, the projection of $c_\sigma$ to
$V_0$ is independent of the lift of $h_\sigma$ and satisfies
\[
  c_0(hk)=c_0(h)+h c_0(k).
\]
Because a reductive algebraic group over $\C$ is linearly reductive, the
associated extension of the trivial representation by $V_0$ splits.  Hence
there is $v_0\in V_0$ with
\[
  c_0(h)=v_0-hv_0.
\]
\end{proof}

\begin{theorem}[Linearity of the functional relation ideal]
\label{thm:functional-linearity}
For the finite direct-sum moment vector in \eqref{eq:direct-system}, the
kernel of
\begin{equation}\label{eq:functional-map-C}
  \Theta_{\Gamma,\C}:
  R_\Gamma[X_1,\ldots,X_N]\longrightarrow\mathscr O(\C),
  \qquad
  X_i\longmapsto U_i(z),
\end{equation}
is generated by polynomials affine-linear in the variables $X_i$.
More precisely, there are $S\in\operatorname{GL}_N(A_\Gamma)$, an integer
$0\le t\le N$, and $g_1,\ldots,g_t\in R_\Gamma$ such that, with
$\mathbf V=S\mathbf U$,
\begin{equation}\label{eq:split-coordinates}
  V_i=g_i\quad(1\le i\le t),
  \qquad
  V_{t+1},\ldots,V_N
  \text{ are algebraically independent over }K_\Gamma.
\end{equation}
\end{theorem}

\begin{proof}
Put $K=K_\Gamma$.  Let $L_h=K(Y)$ be a Picard--Vessiot field for the
homogeneous system and let $L=K(Y,\mathbf U)$ be the Picard--Vessiot field
of the augmented system.  By \cref{lem:base-change-semisimple},
$H=\operatorname{Gal}(L_h/K)$ is reductive.  Apply
\cref{lem:affine-exact-sequence} and let $W$ be the pure-translation
subspace of the constant solution space.

Choose an $H$-equivariant projector $E$ with image $W$.  Then
\[
  P=YEY^{-1}
\]
is fixed by $H$, belongs to $M_N(K)$, and satisfies $P'=[A,P]$.  By
\cref{thm:endomorphism-normal-form}, $P\in M_N(A_\Gamma)$.  Its image and
kernel in $A_\Gamma^N$ are finitely generated projective modules.  Since
$A_\Gamma$ is a Laurent polynomial ring over a field, Swan's theorem makes
both free \cite{SwanLaurent}.  Choosing bases gives
$S\in\operatorname{GL}_N(A_\Gamma)$ such that
\begin{equation}\label{eq:standard-projector}
  SPS^{-1}=P_0=\operatorname{diag}(0_t,I_{N-t}).
\end{equation}

After a constant change of columns in $Y$, we may also assume that
$E=P_0$.  If $B=SY$, then $BE B^{-1}=P_0=E$, so $B$ commutes with $E$ and
is block diagonal:
\[
  B=\operatorname{diag}(B_0,B_W).
\]
Write $\mathbf V=S\mathbf U=(\mathbf V_0,\mathbf V_W)^{\mathsf T}$.
The translation term on the first block is the cocycle from
\cref{lem:affine-exact-sequence}.  Thus, for some constant vector $v_0$,
\[
  \mathbf a=\mathbf V_0+B_0v_0
\]
is fixed by the augmented Galois group and belongs to $K^t$.  The matrices
$S$ and $S^{-1}$ are finite Laurent sums of exponential functions and are
invertible at zero.  Hence
$B_0v_0=\mathbf a-\mathbf V_0$ is a single-valued meromorphic homogeneous
solution at zero after returning through $S^{-1}$.  By
\cref{lem:meromorphic-exclusion}, it vanishes.  Therefore
\begin{equation}\label{eq:V0-in-K}
  \mathbf V_0\in K^t.
\end{equation}

Let $Q=I-P$.  Equation \eqref{eq:standard-projector} and
\eqref{eq:V0-in-K} imply $Q\mathbf U\in K^N$.  Horizontality of $Q$ gives
\begin{equation}\label{eq:QU-equation}
  (Q\mathbf U)'=A(Q\mathbf U)+Q\mathbf f.
\end{equation}
The forcing term $Q\mathbf f$ has only finitely many torus weights.  The
operator $\mathscr L$ in \cref{lem:meromorphic-exclusion} commutes with the
torus and is injective on $K^N$.  Consequently the torus orbit of
$Q\mathbf U$ has finite-dimensional span.  By
\cref{lem:finite-orbit},
\[
  Q\mathbf U=\sum_{\mu\in T}e^{\mu z}\mathbf r_\mu(z),
  \qquad \mathbf r_\mu\in\C(z)^N.
\]
Comparing weights in \eqref{eq:QU-equation} gives
\begin{equation}\label{eq:rmu-vector}
  \mathbf r_\mu'
  =(C-\mu I)\mathbf r_\mu
   +\frac Dz\mathbf r_\mu
   +\frac{\mathbf d_\mu}{z}.
\end{equation}
Pole-order comparison excludes poles away from zero.  A leading term
$z^{-m}v$, $m\ge1$, at zero would give $-mv=Dv$, impossible by
\eqref{eq:global-D-spectrum}.  Hence every $\mathbf r_\mu$ is polynomial,
so
\begin{equation}\label{eq:QU-in-R}
  Q\mathbf U\in R_\Gamma^N.
\end{equation}
It follows that the first $t$ coordinates of $\mathbf V=S\mathbf U$ are the
functions $g_i\in R_\Gamma$ in \eqref{eq:split-coordinates}.

The pure translations act on the remaining block by
\[
  \mathbf V_W\longmapsto\mathbf V_W+B_Wc,
  \qquad c\in\C^{N-t}.
\]
If a polynomial over $K$ vanished at $\mathbf V_W$, all its translates would
vanish.  Since $B_W$ is invertible, the affine substitution
$X\mapsto\mathbf V_W+B_Wc$ is injective over $L$, and a polynomial vanishing
for every $c\in\C^{N-t}$ is zero.  Thus the coordinates of $\mathbf V_W$
are algebraically independent over $K$.

Finally, $\mathbf X\mapsto S\mathbf X$ is an automorphism of
$R_\Gamma[\mathbf X]$.  In the transformed coordinates the kernel of
\eqref{eq:functional-map-C} is exactly
\[
  (V_1-g_1,\ldots,V_t-g_t),
\]
because $R_\Gamma$ embeds in $\mathscr O(\C)$ and the remaining coordinates
are algebraically independent over its fraction field.  Pulling this ideal
back by $S$ proves the theorem.
\end{proof}

\begin{corollary}[Descent to algebraic coefficients]
\label{cor:algebraic-descent}
Suppose the data are defined over $k$.  Let
\[
  R_{\Gamma,k}=k[z]\otimes_k k[\Gamma].
\]
The kernel of
\[
  R_{\Gamma,k}[X_1,\ldots,X_N]
  \longrightarrow k[[z]],
  \qquad X_i\longmapsto U_i(z),
\]
is generated by finitely many polynomials affine-linear in the $X_i$.
\end{corollary}

\begin{proof}
The Taylor coefficients of all moment and exponential functions lie in
$k$.  The natural map
\[
  \C\otimes_k k[[z]]\longrightarrow\C[[z]]
\]
is injective: in a finite tensor sum, choose the complex scalar coefficients
linearly independent over $k$ and compare Taylor coefficients.  Flatness
therefore identifies the complex kernel with the scalar extension of the
kernel over $k$, and the same holds for their affine-linear parts.  Apply
\cref{thm:functional-linearity} and descend the ideal equality by faithful
flatness.  The ring $R_{\Gamma,k}[\mathbf X]$ is Noetherian, so finitely many
affine-linear elements suffice.
\end{proof}

\section{Specialization and the formal polynomial theorem}
\label{sec:specialization}

Let
\[
  \mathscr E_k=k[(k,+)]
\]
be the group algebra of the additive group of $k$, with basis
$\varepsilon_\lambda$.  Every element has finite support.  At the numerical
level we evaluate $\varepsilon_\lambda$ as $e^\lambda$.

\begin{theorem}[Polynomial value transfer]\label{thm:value-transfer}
Let $\mathbf U=(U_1,\ldots,U_N)^{\mathsf T}$ be a finite direct sum of moment
vectors with algebraic data.  Let
\[
  P\in\mathscr E_k[X_1,\ldots,X_N]
\]
and suppose
\begin{equation}\label{eq:value-relation}
  P\big|_{\varepsilon_\lambda=e^\lambda,\,
          X_i=U_i(1)}=0.
\end{equation}
Then there are finitely many polynomials
\begin{equation}\label{eq:linear-functional-generators}
  L_\nu(z,\mathbf X)
  \in\bigl(k[z]\otimes_k k[\Gamma]\bigr)[\mathbf X],
\end{equation}
for a finitely generated additive subgroup $\Gamma\subset k$, such that
\begin{enumerate}[label=\textup{(\roman*)},leftmargin=2.4em]
\item each $L_\nu$ is affine-linear in the moment variables;
\item its functional image under
$\varepsilon_\lambda\mapsto e^{\lambda z}$ and
$X_i\mapsto U_i(z)$ is identically zero;
\item $P$ belongs to the ideal generated by the specializations
$L_\nu(1,\mathbf X)$ in $k[\Gamma][\mathbf X]$.
\end{enumerate}
\end{theorem}

\begin{proof}
Only finitely many exponential indices occur in $P$.  Let $\Gamma$ be the
additive subgroup generated by them and by all endpoint values in the moment
systems.  Adjoin to $\mathbf U$ the constant function $1$ and the finitely
many exponential coordinates needed for $P$ and for the forcing terms.
Homogenize $P$ with the constant coordinate.  By
\cref{thm:Beukers} at $z=1$, the relation lifts to a homogeneous polynomial
functional relation with coefficients in $k[z]$ and the same total degree.

Replace every monomial in exponential coordinates by the corresponding
single group-algebra basis element, using
\[
  \prod_j e^{\lambda_jz}=e^{(\sum_j\lambda_j)z},
\]
and dehomogenize at the constant coordinate.  We obtain
\[
  \widetilde P\in
  \bigl(k[z]\otimes_k k[\Gamma]\bigr)[\mathbf X]
\]
with zero functional image and specialization $P$ at $z=1$.  By
\cref{cor:algebraic-descent}, the functional kernel is generated inside this
polynomial ring by finitely many affine-linear elements $L_\nu$.  Expressing
$\widetilde P$ in their ideal and evaluating at $z=1$ proves the result.
\end{proof}

\begin{remark}[Algebraic evaluation parameters]
A common parameter $\xi\in k^\times$ is redundant.  Replacing every
polynomial $q$ by $\xi q$ and every exponential index $\lambda$ by
$\xi\lambda$ converts a relation at $z=\xi$ into the statement of
\cref{thm:value-transfer} at $z=1$.
\end{remark}

\begin{proof}[Proof of \cref{thm:main}]
Let $x\in\Acal_{\Aone}^{\mathrm{KZ,cpt}}$ have zero numerical period.  It is
represented by a polynomial in finitely many affine-line symbols and
zero-dimensional exponential symbols.  Apply the canonical twisted
polynomial reduction and the treatment of constant and linear polynomials
from \cite{LongLinearKZ}.  Thus $x$ is represented by a polynomial in
finitely many reduced symbols
\[
  \mathsf P(q,s,\Delta),
  \qquad
  \deg q\ge2,
  \quad
  \deg s\le\deg q-2,
\]
and the symbols $\mathsf E(\lambda)$.  Each reduced period is a constant
$k$-linear combination of coordinates in a moment vector.

Apply \cref{thm:value-transfer}.  The numerical polynomial relation lies in
the ideal generated by specializations at $z=1$ of affine-linear functional
relations.  A term
$\varepsilon_\mu X_i$ specializes to $e^\mu U_i(1)$, which is the affine-line
period obtained by replacing the corresponding polynomial $q$ by
$q+\mu$.  Hence every specialized generator is a numerical linear relation
among compact affine-line periods and zero-dimensional exponential periods.
By the linear Kontsevich--Zagier theorem of \cite{LongLinearKZ}, its class in
$\Pcal_{\Aone}^{\mathrm{KZ,cpt}}$ is zero.  Its arbitrary polynomial
multiples are therefore zero in the symmetric-algebra quotient
$\Acal_{\Aone}^{\mathrm{KZ,cpt}}$.  Thus $x=0$, proving injectivity.

\end{proof}

\begin{corollary}[Complete polynomial relation criterion]
For finitely many compact affine-line exponential periods
$P_1,\ldots,P_m$ and every polynomial
$Q\in k[X_1,\ldots,X_m]$, one has
\[
  Q(P_1,\ldots,P_m)=0
\]
if and only if the relation is generated by the linear formal relations of
\cite{LongLinearKZ}, multiplication, and
\eqref{eq:E-zero}--\eqref{eq:constant-shift}.
\end{corollary}

\begin{corollary}[Tensor-separable higher-dimensional periods]
\label{cor:tensor-separable}
Consider finitely many product integrals
\[
  J_a=
  \int_{\Gamma_{a,1}\times\cdots\times\Gamma_{a,d_a}}
  \exp\!\left(\sum_{j=1}^{d_a}q_{a,j}(x_j)\right)
  \prod_{j=1}^{d_a}r_{a,j}(x_j)\,dx_j,
\]
where the $q_{a,j},r_{a,j}$ are algebraic polynomials and every factor chain
is compact with algebraic boundary.  Every polynomial numerical relation
among the $J_a$ is generated, after Fubini, by the one-dimensional formal
relations in \cref{thm:main}.
\end{corollary}

\begin{proof}
Fubini writes each $J_a$ as a product of compact affine-line exponential
periods.  Substitute those products and apply \cref{thm:main}.
\end{proof}

\section{Scope and relation to existing work}\label{sec:scope}

The theorem concerns the tensor-generated algebra of compact exponential
periods on the affine line.  It allows arbitrary polynomial $q$, arbitrary
polynomial differential $r(x)\,dx$, arbitrary degree-zero algebraic
zero-cycles, and polynomial relations of arbitrary degree.  Its content is
not that Beukers' theorem lifts polynomial relations: that lifting is
classical.  Nor is it the effective computation of the relation ideal of a
fixed tuple, which is available in \cite{FischlerRivoalEffective}.  The new
structural assertion is that, uniformly for this family and after adjoining
the required exponential shifts, the entire functional ideal is generated by
affine-linear relations; \cite{LongLinearKZ} then identifies those relations
formally.

The restriction to the affine line and compact chains is essential to the
present arithmetic argument.  Taylor coefficients of \eqref{eq:I} are
algebraic because polynomial differentials on $\Aone$ have polynomial
primitives.  On a general algebraic curve they are ordinarily periods, so
Beukers' theorem no longer applies directly.  Genuinely coupled
higher-dimensional exponential periods are also outside the result.  The
corollary above treats only those obtained from products of the one-dimensional
objects.  Setting the exponential function equal to one in arbitrary
dimension already recovers the ordinary Kontsevich--Zagier conjecture.
Noncompact rapid-decay chains introduce the arithmetic of Gamma values and
require irregular specialization beyond the present framework.

\section*{Acknowledgments and AI-assisted research disclosure}

This manuscript is a companion to \cite{LongEIT,LongLinearKZ}.  During its
preparation the author used OpenAI ChatGPT 5.6 Pro and Anthropic Claude Opus
5 for proof exploration, adversarial checking, symbolic verification,
bibliographic verification, notational consistency, and language editing.
The AI systems are not authors.  The authors reviewed the arguments and
bear full responsibility for every statement, proof, citation, and remaining
error.

\small
\begin{thebibliography}{99}

\bibitem{AdamczewskiRivoal}
B.~Adamczewski and T.~Rivoal,
\emph{Exceptional values of $E$-functions at algebraic points},
Bull. Lond. Math. Soc. \textbf{50} (2018), no.~4, 697--708.
\href{https://doi.org/10.1112/blms.12168}{doi:10.1112/blms.12168}.

\bibitem{Beukers}
F.~Beukers,
\emph{A refined version of the Siegel--Shidlovskii theorem},
Ann. of Math. (2) \textbf{163} (2006), no.~1, 369--379.
\href{https://doi.org/10.4007/annals.2006.163.369}{doi:10.4007/annals.2006.163.369}.

\bibitem{BostanRivoalSalvy}
A.~Bostan, T.~Rivoal, and B.~Salvy,
\emph{Minimization of differential equations and algebraic values of
$E$-functions},
Math. Comp. \textbf{93} (2024), no.~347, 1427--1472.
\href{https://doi.org/10.1090/mcom/3912}{doi:10.1090/mcom/3912}.

\bibitem{CommelinHabeggerHuber}
J.~Commelin, P.~Habegger, and A.~Huber,
\emph{Exponential periods and o-minimality},
arXiv:2007.08280v4 (2025), to appear in
Ann. Inst. Fourier (Grenoble).

\bibitem{CressonViuSos}
J.~Cresson and J.~Viu-Sos,
\emph{On the equality of periods of Kontsevich--Zagier},
J. Th\'eor. Nombres Bordeaux \textbf{34} (2022), no.~2, 323--343.
\href{https://doi.org/10.5802/jtnb.1204}{doi:10.5802/jtnb.1204}.

\bibitem{FischlerRivoalEffective}
S.~Fischler and T.~Rivoal,
\emph{Effective algebraic independence of values of $E$-functions},
Math. Z. \textbf{305} (2023), article~48.
\href{https://doi.org/10.1007/s00209-023-03373-9}{doi:10.1007/s00209-023-03373-9}.

\bibitem{FresanJossen}
J.~Fres\'an and P.~Jossen,
\emph{A non-hypergeometric $E$-function},
Ann. of Math. (2) \textbf{194} (2021), no.~3, 903--942.
\href{https://doi.org/10.4007/annals.2021.194.3.7}{doi:10.4007/annals.2021.194.3.7}.

\bibitem{HuberMullerStach}
A.~Huber and S.~M\"uller-Stach,
\emph{Periods and Nori Motives},
Ergebnisse der Mathematik und ihrer Grenzgebiete, 3.~Folge, vol.~65,
Springer, Cham, 2017.
\href{https://doi.org/10.1007/978-3-319-50926-6}{doi:10.1007/978-3-319-50926-6}.

\bibitem{KontsevichZagier}
M.~Kontsevich and D.~Zagier,
\emph{Periods},
in \emph{Mathematics Unlimited--2001 and Beyond},
Springer, Berlin, 2001, pp.~771--808.
\href{https://doi.org/10.1007/978-3-642-56478-9_39}{doi:10.1007/978-3-642-56478-9\_39}.

\bibitem{LongEIT}
C.~D.~Long and A.-A.~Le~Blanc,
\emph{Algebraic Exponential Integrals: Transcendence and Linear Relations},
companion preprint, August 2026.

\bibitem{LongLinearKZ}
C.~D.~Long and A.-A.~Le~Blanc,
\emph{A Linear Kontsevich--Zagier Theorem for Compact Exponential Periods on
the Affine Line},
companion preprint, August 2026.

\bibitem{SwanLaurent}
R.~G.~Swan,
\emph{Projective modules over Laurent polynomial rings},
Trans. Amer. Math. Soc. \textbf{237} (1978), 111--120.
\href{https://doi.org/10.1090/S0002-9947-1978-0469906-4}{doi:10.1090/S0002-9947-1978-0469906-4}.

\bibitem{vanDerPutSinger}
M.~van~der~Put and M.~F.~Singer,
\emph{Galois Theory of Linear Differential Equations},
Grundlehren der mathematischen Wissenschaften, vol.~328,
Springer, Berlin, 2003.
\href{https://doi.org/10.1007/978-3-642-55750-7}{doi:10.1007/978-3-642-55750-7}.

\end{thebibliography}
\end{document}
